% ============================================================
% fig_ejemplo_contexto.tex — Análisis paso a paso de una oración
% Uso: % ============================================================
% fig_ejemplo_contexto.tex — Análisis paso a paso de una oración
% Uso: % ============================================================
% fig_ejemplo_contexto.tex — Análisis paso a paso de una oración
% Uso: % ============================================================
% fig_ejemplo_contexto.tex — Análisis paso a paso de una oración
% Uso: \input{figures/fig_ejemplo_contexto}
% ============================================================
\begin{figure}[H]
\centering
\begin{tikzpicture}[
    word/.style={draw=gray!40, rounded corners, fill=gray!8,
                 minimum width=1.0cm, minimum height=0.6cm, font=\small},
    focus/.style={draw=blue!60, rounded corners, fill=blue!12,
                  minimum width=1.0cm, minimum height=0.6cm,
                  font=\small\bfseries},
    ctx/.style={draw=green!50!black, rounded corners, fill=green!8,
                minimum width=1.0cm, minimum height=0.6cm, font=\small},
    box/.style={draw=gray!40, rounded corners, fill=white,
                inner sep=6pt, font=\footnotesize, align=left},
    arrow/.style={-{Stealth}, thick, gray},
    step/.style={font=\small\bfseries, blue!70}
]

% ── Oración completa (referencia) ───────────────────────────
\node[font=\small\bfseries, gray] at (3.5, 6.8)
    {Oración: \textit{``Vi a la estrella con el telescopio''}};

% ── PASO 1: token aislado ────────────────────────────────────
\node[step] at (-0.5, 5.9) {Paso 1};
\node[box]  at (4.2, 5.85) {Sin contexto: \textbf{``estrella''} puede ser\\
    \textcolor{blue!70}{$\bullet$ astro celeste} \quad
    \textcolor{orange!80}{$\bullet$ persona famosa}};

\node[focus] at (1.8, 5.85) {estrella};
\node[font=\tiny, red!60] at (1.8, 5.45) {ambigua};

% ── PASO 2: ventana pequeña (±2) ─────────────────────────────
\node[step] at (-0.5, 4.6) {Paso 2};
\node[box]  at (4.5, 4.55) {Ventana $\pm 2$: \textit{``la estrella con el''}\\
    Sigue ambigua: sin información resolutoria.};

\foreach \w/\i/\st in {la/0/ctx, estrella/1/focus, con/2/ctx, el/3/ctx}{
    \pgfmathparse{int(\i==1)}
    \ifnum\pgfmathresult=1
        \node[focus] at (\i*1.1, 4.55) {\w};
    \else
        \node[ctx]   at (\i*1.1, 4.55) {\w};
    \fi
}
\node[font=\tiny, red!60] at (1.1, 4.15) {aún ambigua};

% ── PASO 3: oración completa ────────────────────────────────
\node[step] at (-0.5, 3.2) {Paso 3};
\node[box, fill=green!5] at (5.0, 3.15)
    {Oración completa: \textit{``con el telescopio''}\\
    $\Rightarrow$ instrumento óptico $\Rightarrow$
    \textcolor{green!50!black}{\textbf{astro celeste}}};

\foreach \w/\i in {Vi/0, a/1, la/2, estrella/3, con/4, el/5, telescopio/6}{
    \pgfmathparse{int(\i==3)}
    \ifnum\pgfmathresult=1
        \node[focus] at (\i*1.05, 3.15) {\w};
    \else
        \pgfmathparse{int(\i==6 || \i==5 || \i==4)}
        \ifnum\pgfmathresult=1
            \node[ctx, draw=green!50!black, fill=green!10] at (\i*1.05, 3.15) {\w};
        \else
            \node[word] at (\i*1.05, 3.15) {\w};
        \fi
    \fi
}

% Arco resolutorio
\draw[-{Stealth}, green!60!black, thick]
    (6.3, 2.82) .. controls (6.3, 2.2) and (3.15, 2.2) .. (3.15, 2.82);
\node[font=\tiny, green!50!black] at (4.7, 1.95)
    {\textit{``telescopio''} resuelve la ambigüedad de \textit{``estrella''}};

% ── Conclusión ───────────────────────────────────────────────
\node[draw=blue!25, fill=blue!5, rounded corners, font=\footnotesize,
      inner sep=7pt, align=center] at (3.5, 1.1)
    {El significado correcto requiere acceso a \textbf{toda} la oración,\\
     no solo a una ventana local. Un modelo que no tenga acceso\\
     global al contexto no puede resolver esta ambigüedad.};

\end{tikzpicture}
\caption{Análisis paso a paso de la resolución de ambigüedad en la oración
\textit{``Vi a la estrella con el telescopio''}. En el Paso 1, el token
\textit{estrella} es ambiguo en aislamiento. En el Paso 2, una ventana local
no aporta información resolutoria. Solo en el Paso 3, con acceso a toda la
oración, la palabra \textit{telescopio} permite inferir que \textit{estrella}
se refiere a un astro y no a una persona famosa.}
\label{fig:ejemplo_contexto}
\end{figure}

% ============================================================
\begin{figure}[H]
\centering
\begin{tikzpicture}[
    word/.style={draw=gray!40, rounded corners, fill=gray!8,
                 minimum width=1.0cm, minimum height=0.6cm, font=\small},
    focus/.style={draw=blue!60, rounded corners, fill=blue!12,
                  minimum width=1.0cm, minimum height=0.6cm,
                  font=\small\bfseries},
    ctx/.style={draw=green!50!black, rounded corners, fill=green!8,
                minimum width=1.0cm, minimum height=0.6cm, font=\small},
    box/.style={draw=gray!40, rounded corners, fill=white,
                inner sep=6pt, font=\footnotesize, align=left},
    arrow/.style={-{Stealth}, thick, gray},
    step/.style={font=\small\bfseries, blue!70}
]

% ── Oración completa (referencia) ───────────────────────────
\node[font=\small\bfseries, gray] at (3.5, 6.8)
    {Oración: \textit{``Vi a la estrella con el telescopio''}};

% ── PASO 1: token aislado ────────────────────────────────────
\node[step] at (-0.5, 5.9) {Paso 1};
\node[box]  at (4.2, 5.85) {Sin contexto: \textbf{``estrella''} puede ser\\
    \textcolor{blue!70}{$\bullet$ astro celeste} \quad
    \textcolor{orange!80}{$\bullet$ persona famosa}};

\node[focus] at (1.8, 5.85) {estrella};
\node[font=\tiny, red!60] at (1.8, 5.45) {ambigua};

% ── PASO 2: ventana pequeña (±2) ─────────────────────────────
\node[step] at (-0.5, 4.6) {Paso 2};
\node[box]  at (4.5, 4.55) {Ventana $\pm 2$: \textit{``la estrella con el''}\\
    Sigue ambigua: sin información resolutoria.};

\foreach \w/\i/\st in {la/0/ctx, estrella/1/focus, con/2/ctx, el/3/ctx}{
    \pgfmathparse{int(\i==1)}
    \ifnum\pgfmathresult=1
        \node[focus] at (\i*1.1, 4.55) {\w};
    \else
        \node[ctx]   at (\i*1.1, 4.55) {\w};
    \fi
}
\node[font=\tiny, red!60] at (1.1, 4.15) {aún ambigua};

% ── PASO 3: oración completa ────────────────────────────────
\node[step] at (-0.5, 3.2) {Paso 3};
\node[box, fill=green!5] at (5.0, 3.15)
    {Oración completa: \textit{``con el telescopio''}\\
    $\Rightarrow$ instrumento óptico $\Rightarrow$
    \textcolor{green!50!black}{\textbf{astro celeste}}};

\foreach \w/\i in {Vi/0, a/1, la/2, estrella/3, con/4, el/5, telescopio/6}{
    \pgfmathparse{int(\i==3)}
    \ifnum\pgfmathresult=1
        \node[focus] at (\i*1.05, 3.15) {\w};
    \else
        \pgfmathparse{int(\i==6 || \i==5 || \i==4)}
        \ifnum\pgfmathresult=1
            \node[ctx, draw=green!50!black, fill=green!10] at (\i*1.05, 3.15) {\w};
        \else
            \node[word] at (\i*1.05, 3.15) {\w};
        \fi
    \fi
}

% Arco resolutorio
\draw[-{Stealth}, green!60!black, thick]
    (6.3, 2.82) .. controls (6.3, 2.2) and (3.15, 2.2) .. (3.15, 2.82);
\node[font=\tiny, green!50!black] at (4.7, 1.95)
    {\textit{``telescopio''} resuelve la ambigüedad de \textit{``estrella''}};

% ── Conclusión ───────────────────────────────────────────────
\node[draw=blue!25, fill=blue!5, rounded corners, font=\footnotesize,
      inner sep=7pt, align=center] at (3.5, 1.1)
    {El significado correcto requiere acceso a \textbf{toda} la oración,\\
     no solo a una ventana local. Un modelo que no tenga acceso\\
     global al contexto no puede resolver esta ambigüedad.};

\end{tikzpicture}
\caption{Análisis paso a paso de la resolución de ambigüedad en la oración
\textit{``Vi a la estrella con el telescopio''}. En el Paso 1, el token
\textit{estrella} es ambiguo en aislamiento. En el Paso 2, una ventana local
no aporta información resolutoria. Solo en el Paso 3, con acceso a toda la
oración, la palabra \textit{telescopio} permite inferir que \textit{estrella}
se refiere a un astro y no a una persona famosa.}
\label{fig:ejemplo_contexto}
\end{figure}

% ============================================================
\begin{figure}[H]
\centering
\begin{tikzpicture}[
    word/.style={draw=gray!40, rounded corners, fill=gray!8,
                 minimum width=1.0cm, minimum height=0.6cm, font=\small},
    focus/.style={draw=blue!60, rounded corners, fill=blue!12,
                  minimum width=1.0cm, minimum height=0.6cm,
                  font=\small\bfseries},
    ctx/.style={draw=green!50!black, rounded corners, fill=green!8,
                minimum width=1.0cm, minimum height=0.6cm, font=\small},
    box/.style={draw=gray!40, rounded corners, fill=white,
                inner sep=6pt, font=\footnotesize, align=left},
    arrow/.style={-{Stealth}, thick, gray},
    step/.style={font=\small\bfseries, blue!70}
]

% ── Oración completa (referencia) ───────────────────────────
\node[font=\small\bfseries, gray] at (3.5, 6.8)
    {Oración: \textit{``Vi a la estrella con el telescopio''}};

% ── PASO 1: token aislado ────────────────────────────────────
\node[step] at (-0.5, 5.9) {Paso 1};
\node[box]  at (4.2, 5.85) {Sin contexto: \textbf{``estrella''} puede ser\\
    \textcolor{blue!70}{$\bullet$ astro celeste} \quad
    \textcolor{orange!80}{$\bullet$ persona famosa}};

\node[focus] at (1.8, 5.85) {estrella};
\node[font=\tiny, red!60] at (1.8, 5.45) {ambigua};

% ── PASO 2: ventana pequeña (±2) ─────────────────────────────
\node[step] at (-0.5, 4.6) {Paso 2};
\node[box]  at (4.5, 4.55) {Ventana $\pm 2$: \textit{``la estrella con el''}\\
    Sigue ambigua: sin información resolutoria.};

\foreach \w/\i/\st in {la/0/ctx, estrella/1/focus, con/2/ctx, el/3/ctx}{
    \pgfmathparse{int(\i==1)}
    \ifnum\pgfmathresult=1
        \node[focus] at (\i*1.1, 4.55) {\w};
    \else
        \node[ctx]   at (\i*1.1, 4.55) {\w};
    \fi
}
\node[font=\tiny, red!60] at (1.1, 4.15) {aún ambigua};

% ── PASO 3: oración completa ────────────────────────────────
\node[step] at (-0.5, 3.2) {Paso 3};
\node[box, fill=green!5] at (5.0, 3.15)
    {Oración completa: \textit{``con el telescopio''}\\
    $\Rightarrow$ instrumento óptico $\Rightarrow$
    \textcolor{green!50!black}{\textbf{astro celeste}}};

\foreach \w/\i in {Vi/0, a/1, la/2, estrella/3, con/4, el/5, telescopio/6}{
    \pgfmathparse{int(\i==3)}
    \ifnum\pgfmathresult=1
        \node[focus] at (\i*1.05, 3.15) {\w};
    \else
        \pgfmathparse{int(\i==6 || \i==5 || \i==4)}
        \ifnum\pgfmathresult=1
            \node[ctx, draw=green!50!black, fill=green!10] at (\i*1.05, 3.15) {\w};
        \else
            \node[word] at (\i*1.05, 3.15) {\w};
        \fi
    \fi
}

% Arco resolutorio
\draw[-{Stealth}, green!60!black, thick]
    (6.3, 2.82) .. controls (6.3, 2.2) and (3.15, 2.2) .. (3.15, 2.82);
\node[font=\tiny, green!50!black] at (4.7, 1.95)
    {\textit{``telescopio''} resuelve la ambigüedad de \textit{``estrella''}};

% ── Conclusión ───────────────────────────────────────────────
\node[draw=blue!25, fill=blue!5, rounded corners, font=\footnotesize,
      inner sep=7pt, align=center] at (3.5, 1.1)
    {El significado correcto requiere acceso a \textbf{toda} la oración,\\
     no solo a una ventana local. Un modelo que no tenga acceso\\
     global al contexto no puede resolver esta ambigüedad.};

\end{tikzpicture}
\caption{Análisis paso a paso de la resolución de ambigüedad en la oración
\textit{``Vi a la estrella con el telescopio''}. En el Paso 1, el token
\textit{estrella} es ambiguo en aislamiento. En el Paso 2, una ventana local
no aporta información resolutoria. Solo en el Paso 3, con acceso a toda la
oración, la palabra \textit{telescopio} permite inferir que \textit{estrella}
se refiere a un astro y no a una persona famosa.}
\label{fig:ejemplo_contexto}
\end{figure}

% ============================================================
\begin{figure}[H]
\centering
\begin{tikzpicture}[
    word/.style={draw=gray!40, rounded corners, fill=gray!8,
                 minimum width=1.0cm, minimum height=0.6cm, font=\small},
    focus/.style={draw=blue!60, rounded corners, fill=blue!12,
                  minimum width=1.0cm, minimum height=0.6cm,
                  font=\small\bfseries},
    ctx/.style={draw=green!50!black, rounded corners, fill=green!8,
                minimum width=1.0cm, minimum height=0.6cm, font=\small},
    box/.style={draw=gray!40, rounded corners, fill=white,
                inner sep=6pt, font=\footnotesize, align=left},
    arrow/.style={-{Stealth}, thick, gray},
    step/.style={font=\small\bfseries, blue!70}
]

% ── Oración completa (referencia) ───────────────────────────
\node[font=\small\bfseries, gray] at (3.5, 6.8)
    {Oración: \textit{``Vi a la estrella con el telescopio''}};

% ── PASO 1: token aislado ────────────────────────────────────
\node[step] at (-0.5, 5.9) {Paso 1};
\node[box]  at (4.2, 5.85) {Sin contexto: \textbf{``estrella''} puede ser\\
    \textcolor{blue!70}{$\bullet$ astro celeste} \quad
    \textcolor{orange!80}{$\bullet$ persona famosa}};

\node[focus] at (1.8, 5.85) {estrella};
\node[font=\tiny, red!60] at (1.8, 5.45) {ambigua};

% ── PASO 2: ventana pequeña (±2) ─────────────────────────────
\node[step] at (-0.5, 4.6) {Paso 2};
\node[box]  at (4.5, 4.55) {Ventana $\pm 2$: \textit{``la estrella con el''}\\
    Sigue ambigua: sin información resolutoria.};

\foreach \w/\i/\st in {la/0/ctx, estrella/1/focus, con/2/ctx, el/3/ctx}{
    \pgfmathparse{int(\i==1)}
    \ifnum\pgfmathresult=1
        \node[focus] at (\i*1.1, 4.55) {\w};
    \else
        \node[ctx]   at (\i*1.1, 4.55) {\w};
    \fi
}
\node[font=\tiny, red!60] at (1.1, 4.15) {aún ambigua};

% ── PASO 3: oración completa ────────────────────────────────
\node[step] at (-0.5, 3.2) {Paso 3};
\node[box, fill=green!5] at (5.0, 3.15)
    {Oración completa: \textit{``con el telescopio''}\\
    $\Rightarrow$ instrumento óptico $\Rightarrow$
    \textcolor{green!50!black}{\textbf{astro celeste}}};

\foreach \w/\i in {Vi/0, a/1, la/2, estrella/3, con/4, el/5, telescopio/6}{
    \pgfmathparse{int(\i==3)}
    \ifnum\pgfmathresult=1
        \node[focus] at (\i*1.05, 3.15) {\w};
    \else
        \pgfmathparse{int(\i==6 || \i==5 || \i==4)}
        \ifnum\pgfmathresult=1
            \node[ctx, draw=green!50!black, fill=green!10] at (\i*1.05, 3.15) {\w};
        \else
            \node[word] at (\i*1.05, 3.15) {\w};
        \fi
    \fi
}

% Arco resolutorio
\draw[-{Stealth}, green!60!black, thick]
    (6.3, 2.82) .. controls (6.3, 2.2) and (3.15, 2.2) .. (3.15, 2.82);
\node[font=\tiny, green!50!black] at (4.7, 1.95)
    {\textit{``telescopio''} resuelve la ambigüedad de \textit{``estrella''}};

% ── Conclusión ───────────────────────────────────────────────
\node[draw=blue!25, fill=blue!5, rounded corners, font=\footnotesize,
      inner sep=7pt, align=center] at (3.5, 1.1)
    {El significado correcto requiere acceso a \textbf{toda} la oración,\\
     no solo a una ventana local. Un modelo que no tenga acceso\\
     global al contexto no puede resolver esta ambigüedad.};

\end{tikzpicture}
\caption{Análisis paso a paso de la resolución de ambigüedad en la oración
\textit{``Vi a la estrella con el telescopio''}. En el Paso 1, el token
\textit{estrella} es ambiguo en aislamiento. En el Paso 2, una ventana local
no aporta información resolutoria. Solo en el Paso 3, con acceso a toda la
oración, la palabra \textit{telescopio} permite inferir que \textit{estrella}
se refiere a un astro y no a una persona famosa.}
\label{fig:ejemplo_contexto}
\end{figure}
